\subsection*{1.2 Random Variables of Mixed Type and Singular Type}

This section provides additional examples of probability distributions that are not discrete nor continuous. We use the term \textit{mixed-type distribution} to refer to a distribution whose cdf can be expressed as a convex combination of a discrete random variable's cdf and a continuous random variable's cdf. A probability distribution is considered \textit{singular} if all the probability is concentrated on a subset with measure zero.

\textbf{Example 1.2.1 (Saturated Gaussian Random Variable)} \\
Consider a saturation function $f(x)$ defined by
\[
f(x)\triangleq
\begin{cases}
-1.5 & \text{if } x<-1.5,\\
x & \text{if } -1.5\le x\le 1.5,\\
1.5 & \text{if } x>1.5.
\end{cases}
\]

Let $X$ be a standard Gaussian random variable. The random variable $Y=f(X)$ is confined to the interval $[-1.5,1.5]$ and is a mixture of discrete and continuous random variables. Due to the two probability masses located at the boundary of $[-1.5,1.5]$, the cumulative distribution function $F_Y(y)=\Pr(Y\le y)$ has jump discontinuities at $y=\pm 1.5$. The random variable $Y$ does not have a probability density function because the cdf is not differentiable at $y=\pm 1.5$. Figure 1.1 illustrates the cdf of $Y$. This is an example of distribution of mixed type.

\textbf{Figure 1.1.} A plot of the cumulative distribution function in Example 1.2.1.

A natural question to ask in Example 1.2.1 is the conditional distribution of $X$ given $Y=f(X)$. Suppose we can observe the value of $Y$ and want to infer the conditional distribution of $X$ based on the information contained in $Y$. However, we cannot answer this question directly using elementary probability theory because both the pdf of $Y$ and the joint pdf of $X$ and $Y$ do not exist; the usual formula for conditional pdf $f_{XY}(x,y)/f_Y(y)$ does not apply. Therefore, we must resort to more advanced theory to derive the conditional distribution of $X$ given $Y$.

\textbf{Example 1.2.2 (Dirichlet Distribution)} \\
Consider a positive constant $\beta$ and $n$ positive constants $\alpha_1,\dots,\alpha_n$. For $i=1,2,\dots,n$, let $X_i$ be independent Gamma distributed random variables with shape parameter $\alpha_i$ and scale parameter $\beta$, which we denote by $\Gamma(\alpha_i,\beta)$. The pdf of $X_i$ is given by (1.2).

Define
\[
V\triangleq X_1+X_2+\cdots +X_n
\]
as the sum of these Gamma random variables. The components of the random vector
\[
\mathbf{Y}=(Y_1,Y_2,\dots,Y_n)\triangleq (X_1/V,X_2/V,\dots,X_n/V)
\]
are distributed according to the Dirichlet distribution with parameters $\alpha_1,\dots,\alpha_n$. The random vector $\mathbf{Y}$ lies in the region defined by $y_1+y_2+\cdots +y_n=1$ and $y_i\ge 0$ for all $i$ with probability $1$. The Dirichlet distribution is singular because this region has zero volume in $\mathbb{R}^n$. A sample scatter plot is shown in Fig. 1.2.

\textbf{Figure 1.2.} A scatter plot of Dirichlet distribution in Example 1.2.2, with parameters $\alpha_1=\alpha_3=1$ and $\alpha_2=2$. All sample points are on the plane $x+y+z=1$.

The Dirichlet distribution plays a prominent role in the method of latent Dirichlet allocation, a popular technique in natural language processing. Although the Dirichlet distribution does not have a pdf, we can project the random variables to an $(n-1)$-dimensional subspace and describe the probability distribution in the lower-dimensional space. Since the $n$ random values must sum to $1$, we can consider the first $n-1$ components only. The pdf of $Y_1,\dots,Y_{n-1}$ is given by
\begin{equation}
f(y_1,\dots,y_{n-1})
=
\frac{\Gamma(\alpha_1+\cdots +\alpha_n)}{\prod_{k=1}^{n}\Gamma(\alpha_k)}
\left(\prod_{k=1}^{n-1} y_k^{\alpha_k-1}\right)
(1-y_1-\cdots -y_{n-1})^{\alpha_n-1},
\tag{1.4}
\end{equation}
for $(y_1,y_2,\dots,y_{n-1})$ with $y_1+y_2+\cdots +y_{n-1}\le 1$ and $y_k\ge 0$ for $k=1,\dots,n-1$. We can compute probabilities pertaining to this distribution using this lower-dimensional pdf. When $n=2$, the pdf of $Y_1$ reduces to a Beta distribution,
\[
f(y)=
\frac{\Gamma(\alpha_1+\alpha_2)}{\Gamma(\alpha_1)\Gamma(\alpha_2)}
y^{\alpha_1-1}(1-y)^{\alpha_2-1}
\]
for $0\le y\le 1$.

In Python, we can use the \texttt{dirichlet} function in the \texttt{numpy.random} module to generate a Dirichlet-distributed random vector. The following is an example of drawing a number of samples from a Dirichlet distribution with parameters $(1,1,2)$ using the default random number generator.
\begin{verbatim}
from numpy import random
rng = random.default_rng()      # default random number generator
X = rng.dirichlet((1, 1, 2), 8) # draw 8 samples
print(X)                        # print the random samples as an array
\end{verbatim}

A sample run of this program yields the following output:
\begin{verbatim}
[[0.1844921  0.51764633 0.29786156]
 [0.1943196  0.18637496 0.61930544]
 [0.06687522 0.28048234 0.65264245]
 [0.50722478 0.14452905 0.34824617]
 [0.22316573 0.05919519 0.71763908]
 [0.28439455 0.22016395 0.4954415 ]
 [0.18963489 0.13846004 0.67190507]
 [0.20121907 0.28116356 0.51761736]]
\end{verbatim}

Each row of this array represents a sample from the Dirichlet distribution. Note that the sum of the numbers in each row is equal to $1$, as required by the Dirichlet distribution. In this example, we observe that the third component of each sample is often the largest, due to its high weight in the parameter vector $(1,1,2)$.

In each of the previous examples, there are sets with zero length or area that occur with strictly positive probability. As a result, the probability distributions associated with these examples cannot be represented by probability density functions.

We formally define a singular random variable as follows.

\begin{defbox}{1.1}
A real-valued random variable $X$ is said to be \textit{singular} if there exists a set $S$ with length $0$ such that $X$ takes a value in $S$ with probability $1$. Similarly, a random vector $\mathbf{X}$ with values in $\mathbb{R}^n$ is called \textit{singular} if there exists a set $S$ with zero volume such that $\Pr(\mathbf{X}\in S)=1$.

The Cantor distribution is another example of singular random variables.
\end{defbox}

\textbf{Example 1.2.3 (Cantor Distribution)} \\
Consider the infinite series
\[
U=\sum_{i=1}^{\infty} \frac{R_i}{3^i},
\]
where $R_i$ for $i=1,2,3,\dots$ are independent discrete random variables that take value $0$ or $2$, each with probability $1/2$. A realization of $U$ is a number between $0$ and $1$. If we expand $U$ as a $3$-ary number, the digits are all equal to $0$ or $2$, with probability $1$.

If the first digit $R_1$ is zero, the value of $U$ is restricted to the interval $[0,1/3]$, because the largest possible value given $R_1=0$ is $0.02222\ldots$ in base $3$, which is equal to $1/3$. If $R_1=2$, then the value of $U$ is larger than $2/3$. We thus have
\[
\Pr(0\le U\le 1/3)=\Pr(2/3\le U\le 1)=1/2,\qquad \text{and} \qquad \Pr(1/3<U<2/3)=0.
\]

The cdf $F_U(u)\triangleq \Pr(U\le u)$ is equal to $1/2$ for $u\in (1/3,2/3)$.

By repeating the argument in the previous paragraph and considering the event $R_1=0$ and $R_2=0$ and the event $R_1=0$ and $R_2=2$, we can show that $F_U(u)$ is equal to $1/4$ for $u\in (1/9,2/9)$. Similar analysis shows that $F_U(u)$ is equal to $3/4$ for $u\in (7/9,8/9)$. The cdf of $U$ is flat on the intervals
\[
(1/3,2/3),\ (1/9,2/9),\ (7/9,8/9),
\]
\[
(1/27,2/27),\ (7/27,8/27),\ (19/27,20/27),\ (25/27,26/27),\ldots
\]

The lengths of these intervals sum to $1$. Let $T$ be the union of all these intervals. Then, we have $\Pr(U\in T)=0$. The set $C\triangleq [0,1]\setminus T$ is called the \textit{Cantor set}. This set has zero length, and we have $\Pr(U\in C)=1$.

One can prove that the Cantor set $C$ is uncountable. In fact, each point in the Cantor set is uniquely associated with a $3$-ary number such as $0.0202202\ldots$. If we divide this number by $2$, we obtain an infinite binary sequence. Conversely, an infinite binary sequence is associated with a point in the Cantor set. Because the set of all infinite binary sequences is uncountable, so is the Cantor set. The cdf of the Cantor distribution has zero slope almost everywhere, except at the uncountably many points in the Cantor set.

\noindent$\blacktriangleright$ \textbf{Notes on the Definition of Singular Distribution} According to the definition given earlier, a real-valued random variable whose range $S$ is finite or countably infinite is classified as a singular random variable. In some textbooks, a singular distribution is further required to have zero probability at each point in the set $S$. With this additional requirement, a discrete random variable is not regarded as singular, but the Cantor distribution is truly singular.

In the rest of this book, however, we will not require that each point in the set $S$ in Definition 1.1 has zero probability. Therefore, a discrete random variable is regarded as a singular random variable in this book.

The concept of singular distribution allows us to express the distribution of any real-valued random variable as the sum of a continuous part and a singular part. This result is known as the Lebesgue decomposition theorem. In practical applications, the singular part of the decomposition is often a discrete probability distribution concentrated on a set of at most countably many points.
